% GENERATED FILE. Edit canonical lessons, then run npm run generate:books.
% canonical-source-hash: fnv1a64:0188aecf6becada0
% canonical-lessons: PT-C318-festejar, PT-C318-comemorar, PT-C318-comemoracao, PT-C318-brinde, PT-C318-brindar

\chapter{Vamos festejar}
\label{ch:pt-vamos-festejar}
\hlchaptermodality{318}

\begin{chapteropening}
\textbf{By the end of this chapter:} \emph{I can talk about celebrating and commemorating, and about raising a toast.}

Complete the last lesson of chapter 318: I can talk about celebrating and commemorating, and about raising a toast.
\end{chapteropening}

\section[festejar]{festejar — to celebrate (with a party)}

\label{lesson:PT-C318-festejar}



\begin{quote}
Before the new one: say the Portuguese for emigration, then the Portuguese for the emigrant.
\end{quote}

\subsection*{You'll want to know: festejar}
\textbf{festejar} — \textquotedblleft{}to celebrate (with a party)\textquotedblright{}.

\begin{quote}
Festejámos o aniversário dela num restaurante.
\end{quote}

\begin{quote}
\emph{We celebrated her birthday in a restaurant.}
\end{quote}

Livelier than \emph{celebrar}: music, food, friends.

\begin{cousinweb}
\emph{festejar} is built on \emph{festa} \textquotedblleft{}a party\textquotedblright{}, from Latin \emph{fēsta} \textquotedblleft{}feast days\textquotedblright{} (Lewis \& Short); EtymDB lists it beside Occitan \emph{festejar} and French \emph{festoyer}.
\end{cousinweb}

\begin{grammarlens}[title={What you've built}]
One more word for telling, explaining and describing at B1.
\end{grammarlens}

\subsection*{Your turn}
\begin{itemize}
\raggedright
  \item \emph{Say it:} \emph{festejar}
  \item \emph{Say it:} \emph{festejar}, once more
  \item \emph{Say it:} a sentence of your own with \emph{festejar}
  \item \emph{Recall:} say the Portuguese for emigration, then the Portuguese for the emigrant, then say \emph{festejar} again
\end{itemize}

\begin{tcolorbox}[breakable,colback=teal!4,colframe=teal!35!black,title={Before you move on}]
What does festejar mean? (\textquotedblleft{}To celebrate (with a party)\textquotedblright{}.) Say it once more.
\end{tcolorbox}
\section[comemorar]{comemorar — to commemorate, to mark (a date)}

\label{lesson:PT-C318-comemorar}



\begin{quote}
Before the new one: say the Portuguese for the emigrant, then the Portuguese for to celebrate.
\end{quote}

\subsection*{You'll want to know: comemorar}
\textbf{comemorar} — \textquotedblleft{}to commemorate, to mark (a date)\textquotedblright{}.

\begin{quote}
A cidade comemora o fim da guerra todos os anos.
\end{quote}

\begin{quote}
\emph{The city marks the end of the war every year.}
\end{quote}

In Brazil it is also the everyday word for ‘to celebrate’.

\begin{cousinweb}
\emph{comemorar} is a learned borrowing from Latin \emph{commemorāre} \textquotedblleft{}to call to mind\textquotedblright{} (EtymDB), built on \emph{memor} \textquotedblleft{}mindful\textquotedblright{}.
\end{cousinweb}

\begin{grammarlens}[title={What you've built}]
One more word for telling, explaining and describing at B1.
\end{grammarlens}

\subsection*{Your turn}
\begin{itemize}
\raggedright
  \item \emph{Say it:} \emph{comemorar}
  \item \emph{Say it:} \emph{comemorar}, once more
  \item \emph{Say it:} a sentence of your own with \emph{comemorar}
  \item \emph{Recall:} say the Portuguese for the emigrant, then the Portuguese for to celebrate, then say \emph{comemorar} again
\end{itemize}

\begin{tcolorbox}[breakable,colback=teal!4,colframe=teal!35!black,title={Before you move on}]
What does comemorar mean? (\textquotedblleft{}To commemorate, to mark (a date)\textquotedblright{}.) Say it once more.
\end{tcolorbox}
\section[a comemoração]{a comemoração — the commemoration, the celebration}

\label{lesson:PT-C318-comemoracao}



\begin{quote}
Before the new one: say the Portuguese for to celebrate, then the Portuguese for to commemorate.
\end{quote}

\subsection*{You'll want to know: a comemoração}
\textbf{a comemoração} — \textquotedblleft{}the commemoration, the celebration\textquotedblright{}.

\begin{quote}
A comemoração terminou com um concerto.
\end{quote}

\begin{quote}
\emph{The celebration ended with a concert.}
\end{quote}

Built on \emph{comemorar}.

\begin{cousinweb}
\emph{comemoração} is a learned borrowing from Latin \emph{commemorātiō} \textquotedblleft{}a calling to mind\textquotedblright{} (EtymDB), from \emph{commemorāre}.
\end{cousinweb}

\begin{grammarlens}[title={What you've built}]
One more word for telling, explaining and describing at B1.
\end{grammarlens}

\subsection*{Your turn}
\begin{itemize}
\raggedright
  \item \emph{Say it:} \emph{a comemoração}
  \item \emph{Say it:} \emph{a comemoração}, once more
  \item \emph{Say it:} a sentence of your own with \emph{a comemoração}
  \item \emph{Recall:} say the Portuguese for to celebrate, then the Portuguese for to commemorate, then say \emph{a comemoração} again
\end{itemize}

\begin{tcolorbox}[breakable,colback=teal!4,colframe=teal!35!black,title={Before you move on}]
What does a comemoração mean? (\textquotedblleft{}The commemoration, the celebration\textquotedblright{}.) Say it once more.
\end{tcolorbox}
\section[o brinde]{o brinde — the toast (with drinks); a free gift}

\label{lesson:PT-C318-brinde}



\begin{quote}
Before the new one: say the Portuguese for to commemorate, then the Portuguese for the commemoration.
\end{quote}

\subsection*{You'll want to know: o brinde}
\textbf{o brinde} — \textquotedblleft{}the toast (with drinks); a free gift\textquotedblright{}.

\begin{quote}
No fim do jantar, fizemos um brinde aos noivos.
\end{quote}

\begin{quote}
\emph{At the end of dinner we drank a toast to the bride and groom.}
\end{quote}

In shops, \emph{um brinde} is a free gift.

\begin{cousinweb}
\emph{brinde} goes back to a German toast, \emph{ich bring dir’s} \textquotedblleft{}I bring it to you\textquotedblright{} (EtymDB), which reached Portuguese through Spanish and Italian.
\end{cousinweb}

\begin{grammarlens}[title={What you've built}]
One more word for telling, explaining and describing at B1.
\end{grammarlens}

\subsection*{Your turn}
\begin{itemize}
\raggedright
  \item \emph{Say it:} \emph{o brinde}
  \item \emph{Say it:} \emph{o brinde}, once more
  \item \emph{Say it:} a sentence of your own with \emph{o brinde}
  \item \emph{Recall:} say the Portuguese for to commemorate, then the Portuguese for the commemoration, then say \emph{o brinde} again
\end{itemize}

\begin{tcolorbox}[breakable,colback=teal!4,colframe=teal!35!black,title={Before you move on}]
What does o brinde mean? (\textquotedblleft{}The toast (with drinks); a free gift\textquotedblright{}.) Say it once more.
\end{tcolorbox}
\section[brindar]{brindar — to toast, to drink to}

\label{lesson:PT-C318-brindar}



\begin{quote}
Before the new one: say the Portuguese for the commemoration, then the Portuguese for the toast.
\end{quote}

\subsection*{You'll want to know: brindar}
\textbf{brindar} — \textquotedblleft{}to toast, to drink to\textquotedblright{}.

\begin{quote}
Brindámos à saúde da avó.
\end{quote}

\begin{quote}
\emph{We drank to grandmother’s health.}
\end{quote}

\emph{Brindar a} + person or thing.

\begin{cousinweb}
\emph{brindar} is built on \emph{brinde} \textquotedblleft{}a toast\textquotedblright{} (EtymDB), which goes back to a German drinking phrase.
\end{cousinweb}

\begin{grammarlens}[title={What you've built}]
One more word for telling, explaining and describing at B1.
\end{grammarlens}

\subsection*{Your turn}
\begin{itemize}
\raggedright
  \item \emph{Say it:} \emph{brindar}
  \item \emph{Say it:} \emph{brindar}, once more
  \item \emph{Say it:} a sentence of your own with \emph{brindar}
  \item \emph{Recall:} say the Portuguese for the commemoration, then the Portuguese for the toast, then say \emph{brindar} again
\end{itemize}

\begin{tcolorbox}[breakable,colback=teal!4,colframe=teal!35!black,title={Before you move on}]
What does brindar mean? (\textquotedblleft{}To toast, to drink to\textquotedblright{}.) Say it once more.
\end{tcolorbox}
